\documentclass[10pt,aspectratio=169]{beamer}

\usetheme{Madrid}
\usecolortheme{dove}
\setbeamertemplate{navigation symbols}{}
\setbeamertemplate{footline}[frame number]

\usepackage[round,authoryear]{natbib}
\usepackage{booktabs}
\usepackage{array}
\usepackage{amsmath}
\usepackage{xcolor}

\definecolor{accentblue}{RGB}{0,80,160}
\definecolor{accentred}{RGB}{180,30,30}
\definecolor{lightgray}{RGB}{240,240,240}

\title[Functional Load Forecasting]{Functional Forecasting of Turkish Electricity Load Curves}
\author{Zeynep Bozoklu}
\institute{Thesis Proposal}
\date{\today}

\begin{document}

% ============================================================
\begin{frame}
  \titlepage
\end{frame}

% ============================================================
% PART I: ECONOMIC MOTIVATION
% ============================================================

\begin{frame}{The Daily Load Curve: A Portrait of Economic Life}

  \begin{columns}
  \begin{column}{0.54\textwidth}
    \textbf{Why the curve has the shape it has:}
    \vspace{0.3em}
    \begin{itemize}
      \item \textcolor{accentblue}{\textbf{Hours 0--5 (trough):}} Most people asleep.
        Base-load plants (nuclear, hydro) dominate.
        Demand is low and very predictable.
      \item \textcolor{accentred}{\textbf{Hours 6--8 (morning ramp):}} Simultaneous
        awakening. Factories start shifts, HVAC switches on.
        Demand nearly \textbf{doubles in two hours}.
        Timing is sensitive to temperature and day type.
      \item \textbf{Hours 9--16 (plateau):} Industrial and commercial
        activity at full intensity. Predictable from calendar.
      \item \textcolor{accentred}{\textbf{Hours 17--20 (evening surge):}}
        People come home, cook, use AC or heating. Official forecast
        \textbf{most wrong here} (bias: 614--658 MWh/hour).
      \item \textbf{Hours 21--23 (wind-down):} Declining toward the trough.
    \end{itemize}
  \end{column}
  \begin{column}{0.44\textwidth}
    \begin{block}{The forecasting object}
      $Y_d(h)$ = Turkish electricity load on day $d$, hour $h = 0,\ldots,23$

      \vspace{0.5em}
      Not 24 separate quantities ---\\a single daily \textbf{curve}.
    \end{block}
  \end{column}
  \end{columns}
\end{frame}

% ============================================================
\begin{frame}{Who Uses These Forecasts and Why Accuracy Matters}

  \begin{columns}
  \begin{column}{0.48\textwidth}
    \textbf{The Turkish electricity market (post-2013 liberalisation):}
    \vspace{0.3em}
    \begin{itemize}
      \item \textbf{TEİAŞ} (TSO): publishes the official day-ahead
        load forecast used to schedule grid operations and clear
        the day-ahead market (EPIAS). \emph{This is our benchmark.}
      \item \textbf{Generators} bid into the day-ahead market based
        on expected demand. Better demand forecasts $\Rightarrow$
        more strategic bids.
      \item \textbf{Large industrial consumers} face real-time
        imbalance penalties when their consumption deviates from
        contracted amounts. Better forecasts reduce penalties.
      \item \textbf{Unit commitment:} TEİAŞ decides which power plants
        to start up (hours in advance, significant startup costs).
        A better forecast $\Rightarrow$ fewer unnecessary startups.
    \end{itemize}
  \end{column}
  \begin{column}{0.48\textwidth}
    \begin{block}{The systematic underprediction puzzle}
      \centering
      \vspace{0.3em}
      \small
      \begin{tabular}{lrr}
        \toprule
        Year & Off.\ bias & Our bias \\
        \midrule
        2023 & $+344$ MWh & $+176$ MWh \\
        2024 & $+585$ MWh & $+34$ MWh  \\
        2025 & $+472$ MWh & $+98$ MWh  \\
        \midrule
        Avg. & $+467$ MWh & $+103$ MWh \\
        \bottomrule
      \end{tabular}
      \vspace{0.3em}
    \end{block}
    The official forecast \textbf{always underestimates} actual demand.
    The gap worsened in 2024 (hot summer + demand growth).
    Our rolling-origin model adapts continuously and nearly eliminates
    the bias in 2024 ($+34$ MWh).

    \vspace{0.3em}
    \textit{Why does the official underpredict?}
    Likely: model staleness, growing demand, asymmetric institutional
    incentives (over-procurement is visible; real-time reserve costs are not).
  \end{column}
  \end{columns}
\end{frame}

% ============================================================
% PART II: THE FUNCTIONAL APPROACH
% ============================================================

\begin{frame}{From 24 Numbers to a Curve: Why FDA?}

  \begin{columns}
  \begin{column}{0.5\textwidth}
    \textbf{The standard approach:} fit 24 separate forecasting models,
    one per hour. Problems:
    \begin{itemize}
      \item Ignores the shape of the curve --- the morning ramp is a
        \emph{joint} feature of hours 6, 7, and 8, not an independent event.
      \item Cannot borrow information across hours (hour 6 today predicts
        hour 7 today, not just hour 6 tomorrow).
      \item Cannot easily model ``shape features'' like ramp steepness or
        peak timing.
    \end{itemize}

    \vspace{0.5em}
    \textbf{The functional approach:} treat each day as a single
    mathematical object --- a curve $Y_d(\cdot)$.
    \begin{itemize}
      \item Preserves intraday shape information.
      \item Allows differentiation: $DY_d(h) = \frac{d}{dh}Y_d(h)$
        measures how fast demand is \emph{changing} at each hour.
      \item FPCA compresses 24 correlated observations to $K=2$
        uncorrelated scores that carry 97.3\% of curve variation.
    \end{itemize}
  \end{column}
  \begin{column}{0.46\textwidth}
    \textbf{Representation:} B-spline basis ($M=12$ functions)
    \[
      Y_d(h) \approx \sum_{m=1}^{12} c_{dm}\, B_m(h)
    \]
    then FPCA:
    \[
      Y_d(h) = \mu(h) + \sum_{k=1}^{K} \xi_{dk}\,\phi_k(h) + \varepsilon_d(h)
    \]
    \begin{block}{Forecast reconstruction}
      Once we forecast the $K$ scores $\hat{\xi}_d$, the full
      24-hour load curve is recovered as
      \[
        \hat{Y}_d(h) = \hat\mu(h) + \hat\Phi(h)'\hat\xi_d
      \]
    \end{block}
  \end{column}
  \end{columns}
\end{frame}

% ============================================================
\begin{frame}{FPCA: Compressing the Curve to Two Numbers}

  \begin{columns}
  \begin{column}{0.46\textwidth}
    \textbf{Variance decomposition} (full training data):
    \vspace{0.3em}
    \begin{center}
    \small
    \begin{tabular}{rrrr}
      \toprule
      FPC & Load var. & Cum. & Meaning \\
      \midrule
      1 & 93.4\% & 93.4\% & Level \\
      2 &  3.9\% & 97.3\% & Shape \\
      3 &  1.1\% & 98.4\% & --- \\
      4 &  1.0\% & 99.4\% & --- \\
      \bottomrule
    \end{tabular}
    \end{center}
    \vspace{0.3em}
    $K=2$ exceeds the 95\% cumulative variance rule
    while keeping the score vector parsimonious.

    \vspace{0.5em}
    \textbf{Stability check (script 15):}
    We fit FPCA on 2019--2021, 2020--2022, and 2022--2024 separately.
    All pairwise inner products $> 0.993$.
    \emph{The eigenfunctions do not change} --- static FPCA is
    empirically justified.
  \end{column}
  \begin{column}{0.5\textwidth}
    \textbf{Economic meaning of the two components:}

    \vspace{0.5em}
    \textbf{FPC 1 (``Level'', 93.4\%):}
    Its eigenfunction is nearly flat across all 24 hours.
    A high score $\xi_{d1}$ means uniformly high consumption ---
    a busy industrial day, cold winter, hot summer.
    \emph{Think of it as: how much electricity was used today overall.}

    \vspace{0.5em}
    \textbf{FPC 2 (``Shape'', 3.9\%):}
    Its eigenfunction is positive at morning and evening peaks,
    negative at midday.
    A high $\xi_{d2}$ means pronounced morning/evening peaks relative
    to midday --- typical of residential-heavy days (weekends, holidays).
    A low $\xi_{d2}$ means a flatter profile --- typical of summer
    when AC smooths the curve.
    \emph{Think of it as: how ``peaky'' today's curve is.}
  \end{column}
  \end{columns}
\end{frame}

% ============================================================
% PART III: THE FORECASTING MODEL
% ============================================================

\begin{frame}{The Score Forecasting Model}

  After FPCA, the forecasting problem is:
  \textit{given the score history, predict $(\xi_{d1},\xi_{d2})$ for tomorrow.}

  \vspace{0.5em}
  \begin{columns}
  \begin{column}{0.48\textwidth}
    \textbf{FPCA VAR(1,7)} --- core score dynamics:
    \[
      \xi_d = a + A_1\xi_{d-1} + A_7\xi_{d-7} + u_d
    \]
    \begin{itemize}
      \item \textbf{Lag 7:} Monday demand resembles last Monday's.
        Weekly seasonality is the dominant pattern in electricity use.
      \item \textbf{Lag 1:} yesterday's level and shape adjust the
        weekly baseline for recent demand shifts.
      \item Estimated by OLS equation by equation.
    \end{itemize}

    \vspace{0.3em}
    \textbf{Augmented model:}
    \[
      \xi_d = a + A_1\xi_{d-1} + A_7\xi_{d-7} + Bz_d + u_d
    \]
  \end{column}
  \begin{column}{0.48\textwidth}
    \textbf{Feature blocks in $z_d$} (all predetermined at date $d$):
    \vspace{0.2em}
    \begin{itemize}
      \item \textbf{Calendar:} weekend, annual sine/cosine, summer break
      \item \textbf{Holiday:} fixed holidays, Ramadan, Bayram (Eid),
        holiday eves, post-holiday, bridge days
      \item \textbf{Load shape:} yesterday's mean, peak, range,
        rolling 7- and 28-day averages
      \item \textbf{Derivative features:} slope statistics from
        $DY_{d-1}(h) = \tfrac{d}{dh}Y_{d-1}(h)$
        (mean slope, volatility, min/max, morning ramp,
        evening ramp, end-of-day drop)
    \end{itemize}
    \vspace{0.3em}
    \textit{Holiday and derivative features are the two most important
    additions. Without holiday indicators, the model does not beat
    the official forecast.}
  \end{column}
  \end{columns}
\end{frame}

% ============================================================
\begin{frame}{Evaluation Design: No Leakage, No Shortcuts}

  \begin{columns}
  \begin{column}{0.48\textwidth}
    \textbf{Rolling-origin evaluation}
    (walk-forward cross-validation for time series):
    \vspace{0.3em}
    \begin{enumerate}
      \item Test period: Jan 1 2023 -- Dec 31 2025 (1,096 days)
      \item For each test day $d$:
        \begin{itemize}
          \item Estimate FPCA on all data before $d$
          \item Estimate VAR + features on all data before $d$
          \item Produce one-day-ahead forecast for $d$
          \item Record error
        \end{itemize}
      \item Total: 26,304 hourly forecast errors evaluated
    \end{enumerate}
    \vspace{0.3em}
    At every step, \textbf{only information dated $< d$ is used}.
    This mimics real operational conditions.
  \end{column}
  \begin{column}{0.48\textwidth}
    \textbf{Metrics:}
    \[
      \text{MAE} = \frac{1}{N}\sum_{d,h}|e_{dh}|, \quad
      \text{MAPE} = \frac{1}{N}\sum_{d,h}\frac{|e_{dh}|}{Y_d(h)}
    \]
    \begin{itemize}
      \item \textbf{MAE:} absolute error in MWh; directly interpretable
      \item \textbf{MAPE:} percentage error; scale-free; standard in
        the electricity forecasting literature
      \item \textbf{RMSE:} penalises large errors more
    \end{itemize}

    \vspace{0.3em}
    \textbf{Statistical testing:} Diebold-Mariano test
    ($H_1$: challenger has lower expected absolute error than official).
    Applied pooled (all 26,304 observations) and
    hour by hour (24 separate tests).

    \vspace{0.3em}
    \textit{Comparing MAE numbers without a significance test is not enough.}
  \end{column}
  \end{columns}
\end{frame}

% ============================================================
% PART IV: RESULTS
% ============================================================

\begin{frame}{Main Results}

  \begin{center}
  \small
  \begin{tabular}{p{0.38\textwidth}rrrrc}
    \toprule
    Model & MAE & MAPE & MAE gain & MAPE gain & DM \\
    \midrule
    Seasonal naive $Y_{d-7}$ & 1912 & 5.13\% & $-57.6\%$ & $-62.9\%$ & --- \\
    Official forecast          & 1214 & 3.15\% & 0\%       & 0\%       & ref. \\
    \midrule
    FPCA VAR(1,7) scores       & 1443 & 3.76\% & $-18.9\%$ & $-19.4\%$ & $\times$ \\
    $+$ ramp/calendar          & 1238 & 3.26\% & $-2.1\%$  & $-3.7\%$  & $\times$ \\
    $+$ holiday/load-shape     & 1001 & 2.63\% & $+17.5\%$ & $+16.5\%$ & *** \\
    $+$ derivative features    & \textbf{980}  & \textbf{2.56\%} &
                                 $+19.2\%$ & $+18.7\%$ & *** \\
    \midrule
    $+$ residual step          & 653 & 1.70\% & $+46.2\%$ & $+46.0\%$ & *** \\
    \bottomrule
  \end{tabular}
  \end{center}

  \vspace{0.3em}
  \footnotesize
  DM: Diebold-Mariano test vs.\ official forecast (pooled hourly, $H_1$: challenger better).
  *** $p < 0.001$; $\times$ = not significant or worse than official.

  \vspace{0.3em}
  \begin{itemize}
    \item \textbf{Key insight:} the pure VAR dynamics are not enough.
      Holiday and load-shape features are what make the model competitive.
    \item \textbf{MAPE and MAE tell identical stories} --- the gains are
      not a scale artifact.
    \item Main model beats official by \textbf{19\% MAE} and \textbf{19\% MAPE},
      with $p \approx 10^{-163}$.
  \end{itemize}
\end{frame}

% ============================================================
\begin{frame}{Where the Model Wins and Where It Loses}

  \begin{columns}
  \begin{column}{0.52\textwidth}
    \textbf{MAPE gain by hour (main model vs.\ official):}
    \vspace{0.2em}
    \begin{center}
    \small
    \begin{tabular}{clr}
      \toprule
      Hours & Regime & MAPE gain \\
      \midrule
      0--5  & Overnight trough & $+14\%$ to $+60\%$ \\
      \textcolor{accentred}{6--8}  & \textcolor{accentred}{Morning ramp} & \textcolor{accentred}{$-6\%$ to $-24\%$} \\
      9     & Ramp end & $\approx 0\%$ \\
      10--16 & Daytime plateau & $+20\%$ to $+35\%$ \\
      17     & Evening transition & $+7\%$ \\
      \textcolor{accentred}{18--19} & \textcolor{accentred}{Evening surge} & \textcolor{accentred}{$-7\%$ to $-18\%$} \\
      20--23 & Night wind-down & $+10\%$ to $+20\%$ \\
      \bottomrule
    \end{tabular}
    \end{center}
  \end{column}
  \begin{column}{0.44\textwidth}
    \begin{alertblock}{The ramp-hour puzzle}
      Our model \textbf{loses} at hours 6--8 (morning ramp) and
      18--19 (evening surge) --- the steepest transition hours.

      \vspace{0.3em}
      \textbf{Why?} The official forecast incorporates weather
      (temperature) forecasts that determine the \emph{timing}
      of ramps. Our model uses only historical load patterns.

      \vspace{0.3em}
      A one-hour error in ramp timing propagates across
      the entire ramp --- the slope error is not local.
    \end{alertblock}
    \vspace{0.4em}
    DM test: model beats official at \textbf{19 of 24 hours},
    statistically significant at \textbf{18 of 24 hours} (5\% level).
    The residual step wins at \textbf{24 of 24 hours}.
  \end{column}
  \end{columns}
\end{frame}

% ============================================================
\begin{frame}{Diagnostics: Why the Residual Step Is Motivated}

  \begin{columns}
  \begin{column}{0.52\textwidth}
    \textbf{ACF of forecast errors (Ljung-Box, lag 14):}
    \vspace{0.2em}
    \begin{center}
    \small
    \begin{tabular}{lcc}
      \toprule
      Series & LB statistic & Conclusion \\
      \midrule
      Hour 0 errors & 3,124 & Autocorrelated \\
      Hour 6 errors & 7,227 & Strongly autocorrelated \\
      Hour 8 errors & 994   & Autocorrelated \\
      All 24 hours  & ---   & All reject $H_0$ \\
      \bottomrule
    \end{tabular}
    \end{center}
    \vspace{0.3em}
    \textbf{Interpretation:} errors at ramp hours (especially hour 6)
    are correlated across consecutive days.
    If the model misses the ramp timing on Monday, it tends to miss
    it on Tuesday too.
    \emph{This is a predictable pattern.}
  \end{column}
  \begin{column}{0.44\textwidth}
    \begin{block}{The residual correction}
      For target day $d$, using only past data:
      \[
        \hat{e}_s(h) = Y_s(h) - \hat{Y}_s^{base}(h), \quad s < d
      \]
      Estimate a correction curve $\hat{r}_d(h)$ from recent
      residuals and apply:
      \[
        \hat{Y}_d^{res}(h) = \hat{Y}_d^{base}(h) + \lambda_d\,\hat{r}_d(h)
      \]
      $\lambda_d$ is chosen by rolling-origin cross-validation
      on pre-$d$ data only. \textbf{No leakage.}
    \end{block}
    \vspace{0.3em}
    \textbf{Result:} 46\% MAE gain, 46\% MAPE gain vs.\ official.
    Wins at all 24 of 24 hours including the 5 ramp hours
    where the base model lost.
  \end{column}
  \end{columns}
\end{frame}

% ============================================================
% PART V: VALIDATION
% ============================================================

\begin{frame}{Validation and Robustness}

  \begin{columns}
  \begin{column}{0.48\textwidth}
    \textbf{1. Diebold-Mariano test (pooled):}
    \begin{center}
    \scriptsize
    \begin{tabular}{lr}
      \toprule
      Model & DM $p$-value \\
      \midrule
      FPCA VAR(1,7) only & $\approx 1$ (worse) \\
      $+$ ramp/calendar & $0.996$ (worse) \\
      $+$ holiday/load-shape & $10^{-135}$ \\
      $+$ derivative features & $10^{-163}$ \\
      $+$ residual step & $\approx 0$ \\
      \bottomrule
    \end{tabular}
    \end{center}
    The gains are not sampling variation.

    \vspace{0.5em}
    \textbf{2. Eigenfunction stability:}
    \begin{itemize}
      \item Fit FPCA on 2019--21, 2020--22, 2022--24 separately.
      \item Pairwise inner products: all $> 0.993$.
      \item Static FPCA is empirically justified --- no need
        for dynamic FPCA.
    \end{itemize}
  \end{column}
  \begin{column}{0.48\textwidth}
    \textbf{3. $K=4$ robustness:}
    \begin{center}
    \scriptsize
    \begin{tabular}{ccrrr}
      \toprule
      $K$ & Resid. & MAE & MAPE & MAE gain \\
      \midrule
      2 & No  & 980  & 2.56\% & 19.2\% \\
      2 & Yes & 653  & 1.70\% & 46.2\% \\
      4 & No  & 801  & ---    & 34.0\% \\
      4 & Yes & \textbf{640}  & ---    & \textbf{47.3\%} \\
      \bottomrule
    \end{tabular}
    \end{center}
    $K=2$ is the conservative baseline (parsimony + variance rule).
    $K=4$ is a robustness check.

    \vspace{0.5em}
    \textbf{4. Random forest benchmark:}
    \begin{center}
    \scriptsize
    \begin{tabular}{lrrr}
      \toprule
      Model & MAE & MAPE & MAE gain \\
      \midrule
      FPCA-RF (fixed basis) & 1217 & 3.14\% & $-0.3\%$ \\
      FPCA-OLS (main model) & \textbf{980} & \textbf{2.56\%} & $+19.2\%$ \\
      \bottomrule
    \end{tabular}
    \end{center}
    RF ties the official; OLS beats it by 19\%.
    \textbf{The linear score equation is adequate} ---
    no meaningful nonlinearity left on the table.
  \end{column}
  \end{columns}
\end{frame}

% ============================================================
\begin{frame}{Thesis Contribution}

  \begin{columns}
  \begin{column}{0.52\textwidth}
    \textbf{The coherent story in one flow:}
    \vspace{0.3em}
    \begin{enumerate}
      \item Turkish daily electricity load is a \textbf{curve}, not 24
        separate numbers. FDA preserves intraday shape.
      \item FPCA compresses to \textbf{2 scores} (97.3\% variance).
        Eigenfunctions are empirically stable.
      \item Scores follow a \textbf{VAR(1,7) + features}. Holiday
        and load-shape features are essential --- pure dynamics fail.
      \item The model beats the official forecast by
        \textbf{19\% MAE / 19\% MAPE}, $p \approx 10^{-163}$.
      \item It loses at \textbf{ramp hours} (6--8, 18--19) where
        the official has a weather-forecast advantage.
      \item \textbf{Error ACF} shows remaining serial structure.
        The residual step exploits this: 46\% MAE gain, 24/24 hours.
    \end{enumerate}
  \end{column}
  \begin{column}{0.44\textwidth}
    \textbf{Methodological contributions:}
    \begin{itemize}
      \item Functional representation + FPCA score forecasting
        applied to the Turkish grid
      \item Systematic feature engineering from functional derivative
      \item Rigorous rolling-origin design with DM testing
      \item Empirical test of eigenfunction stability
    \end{itemize}
    \vspace{0.5em}
    \textbf{Economic contributions:}
    \begin{itemize}
      \item Document and explain the official forecast's systematic
        underprediction
      \item Identify ramp hours as the forecasting frontier
        (weather vs.\ historical patterns)
      \item Quantify the value of holiday and shape information
        for grid operators and market participants
    \end{itemize}
  \end{column}
  \end{columns}
\end{frame}

% ============================================================
\begin{frame}{Selected References}
  \footnotesize
  \begin{thebibliography}{99}

    \bibitem[Kokoszka and Reimherr(2017)]{kokoszka_reimherr_2017}
    Kokoszka, P. and Reimherr, M. (2017).
    \newblock \textit{Introduction to Functional Data Analysis}. CRC Press.

    \bibitem[Diebold and Mariano(1995)]{dm_1995}
    Diebold, F.X. and Mariano, R.S. (1995).
    \newblock Comparing predictive accuracy.
    \newblock \textit{Journal of Business \& Economic Statistics}, 13(3), 253--263.

    \bibitem[Hyndman and Ullah(2007)]{hyndman_ullah_2007}
    Hyndman, R.J. and Ullah, M.S. (2007).
    \newblock Robust forecasting of mortality and fertility rates.
    \newblock \textit{Computational Statistics \& Data Analysis}, 51(10), 4942--4956.

    \bibitem[Xie and Hong(2016)]{xie_hong_2016}
    Xie, J. and Hong, T. (2016).
    \newblock GEFCom2014 probabilistic electric load forecasting.
    \newblock \textit{International Journal of Forecasting}, 32(3), 1012--1016.

    \bibitem[Harvey et al.(1997)]{hlnb_1997}
    Harvey, D., Leybourne, S. and Newbold, P. (1997).
    \newblock Testing the equality of prediction mean squared errors.
    \newblock \textit{International Journal of Forecasting}, 13(2), 281--291.

  \end{thebibliography}
\end{frame}

\end{document}
